\documentclass[a4paper,12pt]{ctexart}
\usepackage{xcolor,graphicx,amsmath}
\usepackage[top=1.8cm, bottom=1.8cm, left=1.8cm, right=1.8cm]{geometry}
\usepackage{array,hyperref}
\hypersetup{colorlinks=true,linkcolor=blue!70!black}
\linespread{1.35}
\definecolor{defblue}{RGB}{0,80,160}\definecolor{thinkorange}{RGB}{230,120,20}
\definecolor{formulared}{RGB}{180,30,50}\definecolor{funboxbg}{RGB}{245,248,252}
\definecolor{examplebg}{RGB}{255,250,240}\definecolor{practicegreen}{RGB}{40,130,70}
\newenvironment{thinkbox}{\vspace{0.3cm}\begin{center}\begin{minipage}{0.88\textwidth}\colorbox{funboxbg}{\begin{minipage}{0.94\textwidth}\vspace{0.15cm}\textbf{\textcolor{thinkorange}{【 想一想 】}}}{\vspace{0.15cm}\end{minipage}}\end{minipage}\end{center}\vspace{0.3cm}}
\newenvironment{definition}[1]{\vspace{0.2cm}\begin{center}\begin{minipage}{0.88\textwidth}\fbox{\begin{minipage}{0.92\textwidth}\vspace{0.1cm}\textbf{\textcolor{defblue}{【 #1 】}}}{\vspace{0.1cm}\end{minipage}}\end{minipage}\end{center}\vspace{0.2cm}}
\newenvironment{example}{\vspace{0.2cm}\begin{center}\begin{minipage}{0.88\textwidth}\colorbox{examplebg}{\begin{minipage}{0.94\textwidth}\vspace{0.15cm}\textbf{\textcolor{orange!80!black}{【 例题 】}}}{\vspace{0.15cm}\end{minipage}}\end{minipage}\end{center}\vspace{0.2cm}}
\newenvironment{practice}{\vspace{0.3cm}\begin{center}\begin{minipage}{0.88\textwidth}\colorbox{funboxbg}{\begin{minipage}{0.94\textwidth}\vspace{0.15cm}\textbf{\textcolor{practicegreen}{【 练一练 】}}}{\vspace{0.15cm}\end{minipage}}\end{minipage}\end{center}\vspace{0.2cm}}
\newcommand{\formula}[1]{\begin{center}\textcolor{formulared}{\large\boxed{#1}}\end{center}}
\title{\textcolor{blue!70!black}{\Huge\textbf{化学4 · 有机化学与物质结构}}}
\author{}\date{}
\begin{document}\maketitle\thispagestyle{empty}
\section*{\textcolor{blue!70!black}{致同学}}
这是化学系列的最后一本，也是内容最丰富的一本。前半部分是有机化学——从甲烷到DNA，碳原子搭建的生命骨架。后半部分是物质结构——原子内部电子的排列规则、分子的立体形状、晶体的微观排列。两者合在一起，构成了化学的"底层逻辑"。

\newpage

% ============ 有机化学部分 ============
\section{\textcolor{orange!80!black}{第一课\quad 有机化学入门}}

\subsection*{有机物的特征}
含\textbf{碳}的化合物（除CO、CO$_2$、碳酸及碳酸盐外）都是有机化合物。碳原子最外层4个电子，可以形成\textbf{4个共价键}——这是碳能构成几千万种有机物的根本原因。

\subsection*{同系物与同分异构体}
\textbf{同系物}：结构相似、分子组成相差一个或若干个"CH$_2$"原子团的有机物。如CH$_4$（甲烷）和C$_2$H$_6$（乙烷）互为同系物。
\textbf{同分异构体}：分子式相同但\textbf{结构不同}的化合物。如C$_4$H$_{10}$有两种——正丁烷（直链）和异丁烷（支链）。

\subsection*{烷烃的命名}
选最长的碳链为主链、从离支链最近的一端编号、写"几号位-什么基-某烷"。例：(CH$_3$)$_2$CHCH$_2$CH$_3$ $\rightarrow$ \textbf{2-甲基丁烷}。

\begin{example}
写出C$_5$H$_{12}$的所有同分异构体。

\textbf{解：}三种——正戊烷（直链）、异戊烷（2-甲基丁烷）、新戊烷（2,2-二甲基丙烷）。碳原子数相同时，支链越多——沸点越低。
\end{example}

\begin{practice}
1. 同分异构体的分子式\_\_\_\_\_，结构\_\_\_\_\_。\\
2. 写出C$_3$H$_8$的名称。\\
3. （判断）CH$_3$CH$_2$OH和CH$_3$OCH$_3$互为同分异构体。（\quad）\\
4. （判断）同系物的分子组成相差一个或若干个CH$_2$原子团。（\quad）\\
5. 写出2-甲基丁烷的结构简式\_\_\_\_\_。\\
6. （选择）C$_4$H$_{10}$的同分异构体有几种（\quad）\\
\quad A. 2种 \quad B. 3种 \quad C. 4种 \quad D. 5种
\end{practice}

\newpage
\section{\textcolor{orange!80!black}{第二课\quad 烃——只含碳和氢}}

\subsection*{烷烃——全是单键}
通式C$_n$H$_{2n+2}$。代表：甲烷CH$_4$（天然气主要成分）。化学性质稳定，主要发生\textbf{取代反应}：CH$_4$ + Cl$_2$ $\xrightarrow{\text{光}}$ CH$_3$Cl + HCl。

\subsection*{烯烃——有碳碳双键}
通式C$_n$H$_{2n}$。代表：乙烯C$_2$H$_4$。双键中的一条是"不稳定"的$\pi$键——容易断裂发生\textbf{加成反应}：CH$_2$=CH$_2$ + Br$_2$ $\rightarrow$ CH$_2$BrCH$_2$Br（溴水褪色——检验碳碳双键）。乙烯还能发生\textbf{加聚反应}生成聚乙烯（塑料）。

\subsection*{炔烃——有碳碳三键}
通式C$_n$H$_{2n-2}$。代表：乙炔C$_2$H$_2$（电石气，氧炔焰用于焊接切割金属）。

\subsection*{苯——芳香烃的"祖宗"}
C$_6$H$_6$，六元环状结构。苯环中的碳碳键既不是单键也不是双键——是\textbf{介于两者之间的大$\pi$键}。苯的化学性质：易\textbf{取代}（如硝化反应生成硝基苯）、难\textbf{加成}。

\begin{example}
如何鉴别乙烯和乙烷？

\textbf{解：}通入溴水或酸性高锰酸钾溶液——乙烯使其褪色（加成/氧化），乙烷不反应。
\end{example}

\begin{practice}
1. 烷烃通式\_\_\_\_\_；烯烃通式\_\_\_\_\_。\\
2. 写出乙烯与溴的四氯化碳溶液反应的化学方程式。\\
3. （判断）苯能使酸性高锰酸钾溶液褪色。（\quad）\\
4. 甲烷与氯气在光照下发生\_\_\_\_\_反应。\\
5. 乙炔的分子式是\_\_\_\_\_，属于\_\_\_\_\_烃。\\
6. 写出乙烯发生加聚反应生成聚乙烯的化学方程式\_\_\_\_\_。\\
7. （判断）苯分子中的碳碳键是单键和双键交替排列的。（\quad）
\end{practice}

\newpage
\section{\textcolor{orange!80!black}{第三课\quad 烃的衍生物}}

\subsection*{卤代烃}
烃分子中的氢被卤素（F/Cl/Br/I）取代。例：CH$_3$CH$_2$Br（溴乙烷）。卤代烃在NaOH水溶液中发生\textbf{水解反应}（取代）生成醇。

\subsection*{醇与酚}
\textbf{醇}：羟基（—OH）连接在饱和碳上。代表：乙醇C$_2$H$_5$OH（酒精）。乙醇能发生\textbf{消去反应}——浓硫酸、170℃ $\rightarrow$ 乙烯。
\textbf{酚}：羟基直接连在苯环上。代表：苯酚（石炭酸）——弱酸性，遇FeCl$_3$显紫色。

\subsection*{醛与酮}
\textbf{醛}：含有醛基—CHO。代表：甲醛HCHO（福尔马林）、乙醛CH$_3$CHO。醛基可被氧化（银镜反应、与新制Cu(OH)$_2$反应生成砖红色沉淀——检验醛基）。
\textbf{酮}：羰基C=O两端连碳。代表：丙酮CH$_3$COCH$_3$。

\subsection*{羧酸与酯}
\textbf{羧酸}：含有羧基—COOH。代表：乙酸CH$_3$COOH（醋酸）。有弱酸性。
\textbf{酯}：酸和醇脱水生成的产物——\textbf{酯化反应}：CH$_3$COOH + C$_2$H$_5$OH $\rightleftharpoons$ CH$_3$COOC$_2$H$_5$ + H$_2$O（乙酸乙酯，有果香）。酯可以水解回到酸和醇。

\begin{example}
写出乙醇催化氧化生成乙醛的化学方程式。

\textbf{解：}2CH$_3$CH$_2$OH + O$_2$ $\xrightarrow[\Delta]{\text{Cu}}$ 2CH$_3$CHO + 2H$_2$O。Cu作催化剂。
\end{example}

\begin{practice}
1. 乙醇的官能团是\_\_\_\_\_；乙酸的官能团是\_\_\_\_\_。\\
2. （判断）乙醇在浓硫酸140℃发生消去反应生成乙烯。（\quad）（提示：170℃才消去）\\
3. 写出乙酸与乙醇的酯化反应。\\
4. 检验醛基可用\_\_\_\_\_反应或与\_\_\_\_\_反应生成砖红沉淀。\\
5. （判断）乙醇和甲醚（CH$_3$OCH$_3$）互为同分异构体。（\quad）\\
6. 写出乙醇在浓硫酸、170℃条件下消去反应的化学方程式\_\_\_\_\_。\\
7. 甲醛的水溶液俗称\_\_\_\_\_，可用于浸制生物标本。
\end{practice}

\newpage
\section{\textcolor{orange!80!black}{第四课\quad 生物大分子与合成高分子}}

\subsection*{糖类}
\textbf{单糖}：葡萄糖C$_6$H$_{12}$O$_6$（含醛基和多个羟基——既是醛又是醇）、果糖。\\
\textbf{二糖}：蔗糖（葡萄糖+果糖）、麦芽糖。\\
\textbf{多糖}：淀粉（植物的储能物质）、纤维素（植物的结构物质——细胞壁）。淀粉遇碘变\textbf{蓝}。

\subsection*{油脂}
油脂是高级脂肪酸的甘油酯。酸性水解生成甘油+脂肪酸；碱性水解（\textbf{皂化反应}）生成甘油+脂肪酸盐（肥皂的主要成分）。

\subsection*{蛋白质}
蛋白质由\textbf{氨基酸}通过肽键缩合而成。蛋白质具有\textbf{盐析}（加盐沉淀，可逆）和\textbf{变性}（加热、强酸强碱、重金属盐——不可逆）。检验蛋白质：\textbf{双缩脲反应}。

\subsection*{核酸}
DNA（脱氧核糖核酸）——遗传信息的载体，双螺旋结构。RNA（核糖核酸）——参与蛋白质合成。

\subsection*{合成高分子}
\textbf{塑料}：聚乙烯（食品袋）、聚氯乙烯PVC（水管）、聚苯乙烯（泡沫塑料）。\\
\textbf{合成纤维}：尼龙（聚酰胺）、涤纶（聚酯）。\\
\textbf{合成橡胶}：丁苯橡胶、顺丁橡胶。

\begin{example}
为什么高温能使蛋白质"变性"（如煮鸡蛋蛋白凝固）？这对其营养价值有何影响？

\textbf{解：}高温破坏了蛋白质的\textbf{空间结构}（二、三、四级结构被破坏），肽链展开并互相缠绕凝固。变性使蛋白质更容易被消化酶接触——\textbf{更容易消化吸收}。所以煮鸡蛋比生鸡蛋更富营养价值（也更安全）。
\end{example}

\begin{practice}
1. 淀粉遇碘变\_\_\_\_\_色。\\
2. 油脂在碱性条件下水解叫\_\_\_\_\_反应，产物是\_\_\_\_\_和\_\_\_\_\_。\\
3. （判断）蛋白质变性是可逆过程。（\quad）\\
4. 写出葡萄糖的分子式\_\_\_\_\_。\\
5. （判断）纤维素和淀粉互为同分异构体。（\quad）\\
6. 蛋白质在高温下发生\_\_\_\_\_（盐析/变性），该过程\_\_\_\_\_（可逆/不可逆）。\\
7. 葡萄糖分子中含有的官能团是\_\_\_\_\_（填名称）。
\end{practice}

\newpage
% ============ 物质结构部分 ============
\section{\textcolor{orange!80!black}{第五课\quad 原子结构深化}}

\subsection*{核外电子的排布规则}
电子在核外的排布不是随意的——遵循三个原则：
\begin{enumerate}
  \item \textbf{能量最低原理}：电子优先占据能量最低的轨道（1s$\rightarrow$2s$\rightarrow$2p$\rightarrow$3s$\rightarrow$3p$\rightarrow$4s$\rightarrow$3d…注意4s能量低于3d！）。
  \item \textbf{泡利不相容原理}：每个原子轨道最多容纳两个自旋相反的电子。
  \item \textbf{洪特规则}：电子在能量相同的轨道上分布时，尽可能分占不同轨道且自旋相同。
\end{enumerate}

\subsection*{电子排布式}
\textbf{Fe（26号）}：1s$^2$2s$^2$2p$^6$3s$^2$3p$^6$\textbf{3d$^6$4s$^2$}。注意：写电子排布式时4s写在3d之前（填充电子的顺序），失去电子时也从4s开始失。

\subsection*{元素周期律的深层解释}
原子半径：同一周期从左到右——核电荷递增、电子层数不变、原子核对外层电子吸引力增大——半径\textbf{减小}。\\
第一电离能：同一周期从左到右总体\textbf{增大}。\\
电负性：衡量原子吸引电子能力的标度——F的电负性最大（4.0）。

\begin{example}
写出Cr（24号）的电子排布式。注意它的"特殊"。

\textbf{解：}按正常排布应该是3d$^4$4s$^2$，但实验测定是\textbf{3d$^5$4s$^1$}。因为半充满的3d$^5$结构特别稳定——这是一个洪特规则的特例。同族Cu（29号）也是3d$^{10}$4s$^1$。
\end{example}

\begin{practice}
1. 写出O（8号）的电子排布式。\\
2. （填空）同一周期从左到右，原子半径\_\_\_\_\_，第一电离能总体\_\_\_\_\_。\\
3. 4s和3d哪个先填电子？\_\_\_\_\_；失电子时呢？\_\_\_\_\_。\\
4. Fe$^{2+}$的电子排布式是\_\_\_\_\_。\\
5. （判断）同一周期元素从左到右，电负性逐渐减小。（\quad）\\
6. 写出Cl（17号）的电子排布式\_\_\_\_\_。\\
7. 电负性最大的元素是\_\_\_\_\_。
\end{practice}

\newpage
\section{\textcolor{orange!80!black}{第六课\quad 分子结构与杂化轨道}}

\subsection*{共价键的类型}
\textbf{$\sigma$键}："头顶头"重叠——沿键轴方向。单键都是$\sigma$键。\\
\textbf{$\pi$键}："肩并肩"重叠——在键轴两侧。双键 = 1个$\sigma$ + 1个$\pi$；三键 = 1个$\sigma$ + 2个$\pi$。

\subsection*{杂化轨道理论}
原子在形成分子时，能量相近的原子轨道会"混合"成新的\textbf{杂化轨道}：
\begin{itemize}
  \item \textbf{sp杂化}：直线形（180°）。例：BeCl$_2$、CO$_2$、C$_2$H$_2$。
  \item \textbf{sp$^2$杂化}：平面三角形（120°）。例：BF$_3$、C$_2$H$_4$、苯。
  \item \textbf{sp$^3$杂化}：正四面体形（109.5°）。例：CH$_4$、NH$_3$、H$_2$O。
\end{itemize}

判断杂化类型的方法：\textbf{中心原子的$\sigma$键数 + 孤电子对数 = 杂化轨道的数目}。

\subsection*{VSEPR——预测分子形状}
价层电子对尽可能远离彼此。电子对数目决定了分子的基本构型（直线/三角/四面体…），孤对电子的排斥力大于成键电子对——会"挤压"键角。

\textbf{CH$_4$}：4对成键电子 $\rightarrow$ 正四面体，109.5°。\quad
\textbf{NH$_3$}：3对成键 + 1对孤对 $\rightarrow$ 三角锥，约107°。\\
\textbf{H$_2$O}：2对成键 + 2对孤对 $\rightarrow$ V形，约104.5°。

\begin{example}
判断CO$_2$中碳的杂化方式和分子构型。

\textbf{解：}$\sigma$键数=2（两个C=O双键各含一个$\sigma$），孤电子对数=0。$\therefore$ sp杂化，直线形。
\end{example}

\begin{practice}
1. NH$_3$中N的杂化方式是\_\_\_\_\_，分子构型是\_\_\_\_\_。\\
2. （判断）双键 = 1个$\sigma$ + 1个$\pi$。（\quad）\\
3. H$_2$O的键角约\_\_\_\_\_°，比109.5°\_\_\_\_\_（大/小），因为\_\_\_\_\_。\\
4. BF$_3$中B的杂化方式是\_\_\_\_\_。\\
5. （判断）CH$_4$、NH$_3$、H$_2$O的中心原子杂化方式相同。（\quad）\\
6. C$_2$H$_2$中C的杂化方式是\_\_\_\_\_，分子构型是\_\_\_\_\_。\\
7. （选择）下列分子中键角最大的是（\quad）\\
\quad A. CH$_4$ \quad B. NH$_3$ \quad C. H$_2$O \quad D. CO$_2$
\end{practice}

\newpage
\section{\textcolor{orange!80!black}{第七课\quad 分子间作用力与晶体结构}}

\subsection*{分子间作用力（范德华力）}
分子之间也存在微弱的吸引力。分子量越大、分子间接触面积越大，范德华力越强——沸点越高。

\subsection*{氢键——"最强的分子间力"}
当H原子连接在\textbf{F、O、N}上时，这个H会吸引另一分子中的F/O/N——形成\textbf{氢键}。氢键比共价键弱得多，但远强于普通范德华力。

氢键解释了水的异常性质：水的沸点（100℃）远高于同族H$_2$S（-60℃）——因为水分子间有氢键。冰浮在水面上——也是氢键把水分子"撑开"成六角形空旷结构。

\subsection*{四大晶体类型}

\begin{center}
\begin{tabular}{|c|c|c|c|c|}
\hline
\textbf{类型} & \textbf{构成粒子} & \textbf{作用力} & \textbf{性质} & \textbf{实例} \\
\hline
离子晶体 & 阴、阳离子 & 离子键 & 硬而脆、熔沸点高、熔融导电 & NaCl、KNO$_3$\\
分子晶体 & 分子 & 分子间力 & 熔沸点低、不导电 & 干冰、冰、碘\\
共价晶体 & 原子 & 共价键 & 极硬、极高沸点 & 金刚石、SiO$_2$\\
金属晶体 & 金属阳离子+自由电子 & 金属键 & 导电导热、延展 & Cu、Fe、Al\\
\hline
\end{tabular}
\end{center}

\begin{example}
为什么NaCl的熔点（801℃）远高于干冰CO$_2$（-78.5℃升华）？

\textbf{解：}NaCl是\textbf{离子晶体}——Na$^+$和Cl$^-$之间的离子键非常强，破坏需要大量能量。CO$_2$是\textbf{分子晶体}——分子间只有微弱的范德华力，极易破坏。
\end{example}

\begin{thinkbox}
金刚石和石墨都是碳——但金刚石是自然界最硬的物质（共价晶体），石墨则软到能做铅笔芯（层状结构，层间是弱范德华力）。同样的原子，不同的排列——这就是\textbf{同素异形体}的奇妙之处，也是物质结构的核心思想：\textbf{性质由结构决定，不只是组成。}
\end{thinkbox}

\begin{practice}
1. （判断）金刚石和石墨都是碳的同素异形体。（\quad）\\
2. SiO$_2$属于\_\_\_\_\_晶体；干冰CO$_2$属于\_\_\_\_\_晶体。\\
3. 氢键存在于H与\_\_\_\_\_、\_\_\_\_\_、\_\_\_\_\_之间。\\
4. 水的沸点异常高，是因为水分子间存在\_\_\_\_\_。\\
5. （判断）干冰升华时破坏的是共价键。（\quad）\\
6. （选择）下列物质中属于共价晶体的是（\quad）\\
\quad A. NaCl \quad B. CO$_2$ \quad C. 金刚石 \quad D. Cu\\
7. 金属晶体中，自由电子和金属阳离子之间的作用力称为\_\_\_\_\_。
\end{practice}

\vspace{0.5cm}
\begin{center}\rule{0.7\textwidth}{1pt}\vspace{0.4cm}
{\large\textcolor{blue!70!black}{\textbf{—— 化学系列（4本）完 ——}}}
\vspace{0.3cm}\textcolor{gray}{\small 后续系列：生物4本 · 地理2本}
\end{center}

\newpage

\appendix
\section*{\textcolor{blue!70!black}{附录A\quad 元素周期表}}

\begin{figplaceholder}
此处插入完整元素周期表（118元素，彩色，含中文名称与相对原子质量）。\\
本图化学1-4通用（见 figure/requirements/periodic-table.txt）
\end{figplaceholder}

\subsection*{1-20号元素速查表}

\begin{center}
\begin{tabular}{|c|c|c|c|c|}
\hline
\textbf{序数} & \textbf{符号} & \textbf{名称} & \textbf{相对原子质量} & \textbf{常见化合价} \\
\hline
1 & H & 氢 & 1.008 & +1 \\
2 & He & 氦 & 4.003 & 0 \\
3 & Li & 锂 & 6.941 & +1 \\
4 & Be & 铍 & 9.012 & +2 \\
5 & B & 硼 & 10.81 & +3 \\
6 & C & 碳 & 12.01 & +2,+4,-4 \\
7 & N & 氮 & 14.01 & -3,+3,+5 \\
8 & O & 氧 & 16.00 & -2 \\
9 & F & 氟 & 19.00 & -1 \\
10 & Ne & 氖 & 20.18 & 0 \\
11 & Na & 钠 & 22.99 & +1 \\
12 & Mg & 镁 & 24.31 & +2 \\
13 & Al & 铝 & 26.98 & +3 \\
14 & Si & 硅 & 28.09 & +4,-4 \\
15 & P & 磷 & 30.97 & -3,+3,+5 \\
16 & S & 硫 & 32.07 & -2,+4,+6 \\
17 & Cl & 氯 & 35.45 & -1,+1,+5,+7 \\
18 & Ar & 氩 & 39.95 & 0 \\
19 & K & 钾 & 39.10 & +1 \\
20 & Ca & 钙 & 40.08 & +2 \\
\hline
\end{tabular}
\end{center}

\vspace{0.3cm}
\textcolor{gray}{\small 注：相对原子质量为近似值，教学用。稀有气体（He/Ne/Ar）化合价为0。}

\newpage
\section*{\textcolor{defblue}{参考答案}}

\subsection*{第一课}
1. 相同、不同\\
2. 丙烷\\
3. （√）\\
4. （√）\\
5. (CH$_3$)$_2$CHCH$_2$CH$_3$\\
6. A（2种：正丁烷和异丁烷）

\subsection*{第二课}
1. C$_n$H$_{2n+2}$、C$_n$H$_{2n}$\\
2. CH$_2$=CH$_2$ + Br$_2$ $\rightarrow$ CH$_2$BrCH$_2$Br\\
3. （×）\\
4. 取代\\
5. C$_2$H$_2$、炔\\
6. $n$CH$_2$=CH$_2$ $\xrightarrow{\text{催化剂}}$ [CH$_2$—CH$_2$]$_n$\\
7. （×）苯环中碳碳键介于单双键之间，是独特的大$\pi$键

\subsection*{第三课}
1. —OH（羟基）、—COOH（羧基）\\
2. （×）170℃才消去生成乙烯，140℃生成乙醚\\
3. CH$_3$COOH + C$_2$H$_5$OH $\rightleftharpoons$ CH$_3$COOC$_2$H$_5$ + H$_2$O\\
4. 银镜、新制Cu(OH)$_2$\\
5. （√）分子式均为C$_2$H$_6$O，结构不同\\
6. C$_2$H$_5$OH $\xrightarrow[\text{浓H$_2$SO$_4$}]{170℃}$ CH$_2$=CH$_2$$\uparrow$ + H$_2$O\\
7. 福尔马林

\subsection*{第四课}
1. 蓝\\
2. 皂化、甘油、脂肪酸盐（肥皂）\\
3. （×）不可逆\\
4. C$_6$H$_{12}$O$_6$\\
5. （×）n值不同，分子式不同\\
6. 变性、不可逆\\
7. 醛基（—CHO）和羟基（—OH）

\subsection*{第五课}
1. 1s$^2$2s$^2$2p$^4$\\
2. 减小、增大\\
3. 4s、4s\\
4. 1s$^2$2s$^2$2p$^6$3s$^2$3p$^6$3d$^6$（或[Ar]3d$^6$）\\
5. （×）电负性增大\\
6. 1s$^2$2s$^2$2p$^6$3s$^2$3p$^5$（或[Ne]3s$^2$3p$^5$）\\
7. F（氟）

\subsection*{第六课}
1. sp$^3$、三角锥\\
2. （√）\\
3. 104.5、小、孤对电子排斥力大于成键电子对\\
4. sp$^2$\\
5. （√）均为sp$^3$杂化\\
6. sp、直线形\\
7. D（CO$_2$直线形，键角180°）

\subsection*{第七课}
1. （√）\\
2. 共价（原子）、分子\\
3. F、O、N\\
4. 氢键\\
5. （×）破坏的是分子间作用力（范德华力）\\
6. C（金刚石）\\
7. 金属键

\end{document}
